\chapter{編譯問題排查}\label{appendix-a}

使用模板時，若遇到編譯失敗、圖片找不到或 PDF 沒有更新，可以先依本附錄的順序檢查。建議保留終端機中的錯誤訊息，從第一個錯誤指出的檔案與位置開始處理，修正後再重新編譯。

這份檔案也示範附錄的排版方式。換成自己的論文時，可以保留章標題的寫法，將標題與內文替換成問卷、補充數據或其他資料，再到 \code{main.tex} 確認附錄的裝訂順序。

\section{指令執行失敗}

若終端機顯示找不到 \code{xelatex} 或 \code{latexmk}，請先確認已完成安裝，再重新開啟終端機，執行

\begin{lstlisting}[language=bash]
xelatex --version
\end{lstlisting}

確認能夠顯示版本資訊後，再嘗試編譯。若找不到的是 \code{make}，則可以先使用以下指令，直接產生 PDF：

\begin{lstlisting}[language=bash]
texlua scripts/gen.lua
latexmk -xelatex main.tex
\end{lstlisting}

若錯誤訊息指出找不到 \code{main.tex} 或 \code{.tw-thesis-config.tex}，請檢查終端機目前所在的資料夾。編譯指令應在專案根目錄執行，也就是能看到 \code{main.tex}、\code{config.yml} 與 \code{contents/} 的位置；\code{.tw-thesis-config.tex} 則是 \code{scripts/gen.lua} 產生的，先執行過一次就會出現。若目前位於章節資料夾，請先切換回根目錄再試一次。

\section{修改設定後無法編譯}

修改設定後出現錯誤時，請先看訊息指出哪個檔案、哪一個欄位。若包含 \code{unknown key}，可以對照 \code{schools/default.yml} 檢查欄位拼字與所在層級；若是 \code{YAML} 解析錯誤，則先確認縮排使用空白，且冒號後保留空白，再檢查引號是否成對。

加入新設定時，請先找出檔案中是否已有相同區塊，避免在同一層重複建立 \code{layout:}。例如要同時修改字級與左側邊界，應把兩項設定放在同一個 \code{layout:} 下，並讓 \code{left} 再往內縮排一層：

\begin{lstlisting}[language=yaml]
layout:
  font-size: 12
  margin:
    left: 3.5
\end{lstlisting}

若一次修改了多個設定，暫時無法判斷錯誤來源，可以先將剛加入的內容複製到別處保存，再移除這些設定，確認原本的文件能正常編譯。之後逐項加回，每次存檔後檢查結果，較容易找到需要修正的位置。

\section{圖片或引用找不到}

圖片無法載入時，請先確認檔案確實放在 \code{figures/} 的指定位置，再比對路徑中的檔名、副檔名與大小寫。路徑是從專案根目錄算起的，不是從目前章節所在的資料夾開始。若曾經移動或重新命名圖片，也要一併修改章節中的路徑。

圖表引用顯示成 \code{??} 時，請檢查圖表是否已加上 \code{\label}，以及名稱是否與內文的 \code{\ref} 完全相同；也別忘了 \LaTeX{} 需要編譯兩次才會解出交叉引用，\code{latexmk} 會自動處理。若是文獻引用出錯，則到 \code{.bib} 檔案確認識別名稱，再檢查 \code{config.yml} 的 \code{bibliography.path} 是否包含該文獻檔。修正後重新編譯，確認引用已顯示為編號或文獻資料。

\section{PDF 沒有更新}

修改內容後，請先確認終端機是否已完成這次編譯，以及是否出現錯誤。編譯失敗時，資料夾中可能仍保留上一次的 \code{main.pdf}，因此能開啟 PDF 並不代表新內容已經寫入。請先修正錯誤，再重新執行編譯指令。

若使用 \code{make watch}，請確認監看程序仍在執行，並且修改的檔案已經存檔。編譯成功後畫面仍未更新時，可以關閉再重新開啟 PDF，同時確認檢視器開啟的是目前專案中的 \code{main.pdf}，避免看到先前複製到其他位置的版本。若編譯結果卡在舊的狀態，執行 \code{make clean} 清掉暫存檔後再編譯一次，通常就能解決。
